% --- Archon physics formalization source begin ---
% archon:physics
% archon:covers PhyXMiniProblems/problem_phyx_mini_0793.lean
% archon:source-report reports/phyx_mini/problem_phyx_mini_0793.source.json

\chapter{Physics problem phyx\_mini\_0793}
\label{ch:PhyXMiniProblems_problem_phyx_mini_0793}

\paragraph{Problem source.}
\#\# Physical scenario

A motorcyclist heading east through a small town accelerates at a constant \textbackslash{}( 4.0 \textbackslash{}, \textbackslash{}text\{m/s\}\textasciicircum{}2 \textbackslash{}) after he leaves the city limits as shown in figure. At time \textbackslash{}( t = 0 \textbackslash{}) he is \textbackslash{}( 5.0 \textbackslash{}, \textbackslash{}text\{m\} \textbackslash{}) east of the city-limits signpost while he moves east at \textbackslash{}( 15 \textbackslash{}, \textbackslash{}text\{m/s\} \textbackslash{}).

\#\# Question

Where is he when his speed is \textbackslash{}( 25 \textbackslash{}, \textbackslash{}text\{m/s\} \textbackslash{})?

\#\# Answer choices

- A: A: 27m

- B: B: 55m

- C: C: 45m

- D: D: 65m

\#\# Figure caption (auxiliary; use the image as primary evidence)

The image is a diagram illustrating the motion of a motorcycle traveling along a straight path, labeled as the x-axis (east direction). 

1. **Scenes and Objects:**
   - On the left, there's a sign labeled "OSAGE."
   - A motorcyclist is depicted in two positions: at the initial position and later at an unknown later position.

2. **Initial Position:**
   - The motorcyclist at position \textbackslash{}( O \textbackslash{}) has an initial velocity (\textbackslash{}( v\_\{0x\} \textbackslash{})) of 15 m/s, shown by a green arrow pointing to the right.
   - The starting position (\textbackslash{}( x\_0 \textbackslash{})) is 5.0 m away from the origin.
   - The time at this initial position (\textbackslash{}( t \textbackslash{})) is 0 seconds.

3. **Later Position:**
   - The motorcyclist is depicted again at an unknown position (\textbackslash{}( x = ? \textbackslash{})) at a later time (\textbackslash{}( t = 2.0 \textbackslash{}) s).
   - The velocity at this later position is shown with a green arrow and labeled \textbackslash{}( v\_\{x\} = ? \textbackslash{}).

4. **Additional Information:**
   - The acceleration (\textbackslash{}( a\_x \textbackslash{})) is given as 4.0 m/s², represented by an orange arrow along the direction of motion.

This diagram is likely part of a physics problem related to constant acceleration and kinematic equations.

\#\# Dataset metadata

Kinematics; Physical Model Grounding Reasoning, Multi-Formula Reasoning

\paragraph{Recorded answer/context.}
B: B: 55m

\paragraph{Figure/image path.}
/root/proposal\_for\_physic/science-mango/phyx\_mini\_run/phyx\_data/test\_image/793.png

\paragraph{Formalization target.}
create a compiling Lean file with sorry bodies at `PhyXMiniProblems/problem\_phyx\_mini\_0793.lean`. The Lean declarations must preserve the physical quantities, dimensions or dimensional roles, figure labels, governing-law hypotheses, and final relation expressed by this problem.
Use Mathlib/Physlib names found through LeanExplore where available. If a physics API is missing, introduce faithful local abstractions rather than scalar placeholder aliases.

\begin{theorem}[Physics formalization target]
\label{thm:physics:phyx_mini_0793:target}
\lean{PhyXMiniProblems.ProblemPhyXMini0793.problem_phyx_mini_0793}
\uses{def:physics:phyx-mini-0793:phyxminiproblems-problemphyxmini0793-motorcyclemotionsetup, def:physics:phyx-mini-0793:phyxminiproblems-problemphyxmini0793-matchesproblemstatement, def:physics:phyx-mini-0793:phyxminiproblems-problemphyxmini0793-matchessuppliedmotorcyclefigure, def:physics:phyx-mini-0793:phyxminiproblems-problemphyxmini0793-hasphysicaleastwardqueryevent, def:physics:phyx-mini-0793:phyxminiproblems-problemphyxmini0793-satisfiesconstantaccelerationkinematics, def:physics:phyx-mini-0793:phyxminiproblems-problemphyxmini0793-queriedpositioneastofsignpostmeters, def:physics:phyx-mini-0793:phyxminiproblems-problemphyxmini0793-recordedanswerchoice, def:physics:phyx-mini-0793:phyxminiproblems-problemphyxmini0793-matchesdisplayedposition, def:physics:phyx-mini-0793:phyxminiproblems-problemphyxmini0793-isuniquematchingdisplayedchoice}
The assigned autoformalize agent should translate this physics problem into Lean declarations in the covered file, with theorem and lemma proof bodies written as `by sorry`.
\end{theorem}
\begin{proof}
This is an autoformalization task, not a proof task. Produce statements that can later be proved by the physics prover without weakening the physical model.
\end{proof}

% --- Archon declaration topology begin ---
\section{Declaration topology}
\begin{definition}[Signed Position Quantity]
  \label{def:physics:phyx-mini-0793:phyxminiproblems-problemphyxmini0793-signedpositionquantity}
  \lean{PhyXMiniProblems.ProblemPhyXMini0793.SignedPositionQuantity}
  A signed position on the east-positive x-axis
\end{definition}
\begin{proof}
  This is the declared model definition of Signed Position Quantity.
\end{proof}
\begin{definition}[Time Quantity]
  \label{def:physics:phyx-mini-0793:phyxminiproblems-problemphyxmini0793-timequantity}
  \lean{PhyXMiniProblems.ProblemPhyXMini0793.TimeQuantity}
  A physical time coordinate
\end{definition}
\begin{proof}
  This is the declared model definition of Time Quantity.
\end{proof}
\begin{definition}[Signed Velocity Quantity]
  \label{def:physics:phyx-mini-0793:phyxminiproblems-problemphyxmini0793-signedvelocityquantity}
  \lean{PhyXMiniProblems.ProblemPhyXMini0793.SignedVelocityQuantity}
  A signed one-dimensional velocity component
\end{definition}
\begin{proof}
  This is the declared model definition of Signed Velocity Quantity.
\end{proof}
\begin{definition}[Speed Quantity]
  \label{def:physics:phyx-mini-0793:phyxminiproblems-problemphyxmini0793-speedquantity}
  \lean{PhyXMiniProblems.ProblemPhyXMini0793.SpeedQuantity}
  A nonnegative speed magnitude
\end{definition}
\begin{proof}
  This is the declared model definition of Speed Quantity.
\end{proof}
\begin{definition}[Signed Acceleration Quantity]
  \label{def:physics:phyx-mini-0793:phyxminiproblems-problemphyxmini0793-signedaccelerationquantity}
  \lean{PhyXMiniProblems.ProblemPhyXMini0793.SignedAccelerationQuantity}
  A signed one-dimensional acceleration component
\end{definition}
\begin{proof}
  This is the declared model definition of Signed Acceleration Quantity.
\end{proof}
\begin{definition}[position Readout]
  \label{def:physics:phyx-mini-0793:phyxminiproblems-problemphyxmini0793-positionreadout}
  \lean{PhyXMiniProblems.ProblemPhyXMini0793.positionReadout}
  \uses{def:physics:phyx-mini-0793:phyxminiproblems-problemphyxmini0793-signedpositionquantity}
  Read a signed position in the length unit selected by coherent units
\end{definition}
\begin{proof}
  This is the declared model definition of position Readout.
\end{proof}
\begin{definition}[time Readout]
  \label{def:physics:phyx-mini-0793:phyxminiproblems-problemphyxmini0793-timereadout}
  \lean{PhyXMiniProblems.ProblemPhyXMini0793.timeReadout}
  \uses{def:physics:phyx-mini-0793:phyxminiproblems-problemphyxmini0793-timequantity}
  Read a physical time in the time unit selected by coherent units
\end{definition}
\begin{proof}
  This is the declared model definition of time Readout.
\end{proof}
\begin{definition}[velocity Readout]
  \label{def:physics:phyx-mini-0793:phyxminiproblems-problemphyxmini0793-velocityreadout}
  \lean{PhyXMiniProblems.ProblemPhyXMini0793.velocityReadout}
  \uses{def:physics:phyx-mini-0793:phyxminiproblems-problemphyxmini0793-signedvelocityquantity}
  Read a signed velocity in the speed unit induced by coherent units
\end{definition}
\begin{proof}
  This is the declared model definition of velocity Readout.
\end{proof}
\begin{definition}[speed Readout]
  \label{def:physics:phyx-mini-0793:phyxminiproblems-problemphyxmini0793-speedreadout}
  \lean{PhyXMiniProblems.ProblemPhyXMini0793.speedReadout}
  \uses{def:physics:phyx-mini-0793:phyxminiproblems-problemphyxmini0793-speedquantity}
  Read a speed magnitude in the speed unit induced by coherent units
\end{definition}
\begin{proof}
  This is the declared model definition of speed Readout.
\end{proof}
\begin{definition}[acceleration Readout]
  \label{def:physics:phyx-mini-0793:phyxminiproblems-problemphyxmini0793-accelerationreadout}
  \lean{PhyXMiniProblems.ProblemPhyXMini0793.accelerationReadout}
  \uses{def:physics:phyx-mini-0793:phyxminiproblems-problemphyxmini0793-signedaccelerationquantity}
  Read a signed acceleration in the units induced by coherent units
\end{definition}
\begin{proof}
  This is the declared model definition of acceleration Readout.
\end{proof}
\begin{definition}[position In Meters]
  \label{def:physics:phyx-mini-0793:phyxminiproblems-problemphyxmini0793-positioninmeters}
  \lean{PhyXMiniProblems.ProblemPhyXMini0793.positionInMeters}
  \uses{def:physics:phyx-mini-0793:phyxminiproblems-problemphyxmini0793-signedpositionquantity, def:physics:phyx-mini-0793:phyxminiproblems-problemphyxmini0793-positionreadout}
  SI metre readout of a signed position
\end{definition}
\begin{proof}
  This is the declared model definition of position In Meters.
\end{proof}
\begin{definition}[time In Seconds]
  \label{def:physics:phyx-mini-0793:phyxminiproblems-problemphyxmini0793-timeinseconds}
  \lean{PhyXMiniProblems.ProblemPhyXMini0793.timeInSeconds}
  \uses{def:physics:phyx-mini-0793:phyxminiproblems-problemphyxmini0793-timequantity, def:physics:phyx-mini-0793:phyxminiproblems-problemphyxmini0793-timereadout}
  SI second readout of a physical time coordinate
\end{definition}
\begin{proof}
  This is the declared model definition of time In Seconds.
\end{proof}
\begin{definition}[velocity In Meters Per Second]
  \label{def:physics:phyx-mini-0793:phyxminiproblems-problemphyxmini0793-velocityinmeterspersecond}
  \lean{PhyXMiniProblems.ProblemPhyXMini0793.velocityInMetersPerSecond}
  \uses{def:physics:phyx-mini-0793:phyxminiproblems-problemphyxmini0793-signedvelocityquantity, def:physics:phyx-mini-0793:phyxminiproblems-problemphyxmini0793-velocityreadout}
  SI metre-per-second readout of a signed velocity component
\end{definition}
\begin{proof}
  This is the declared model definition of velocity In Meters Per Second.
\end{proof}
\begin{definition}[speed In Meters Per Second]
  \label{def:physics:phyx-mini-0793:phyxminiproblems-problemphyxmini0793-speedinmeterspersecond}
  \lean{PhyXMiniProblems.ProblemPhyXMini0793.speedInMetersPerSecond}
  \uses{def:physics:phyx-mini-0793:phyxminiproblems-problemphyxmini0793-speedquantity, def:physics:phyx-mini-0793:phyxminiproblems-problemphyxmini0793-speedreadout}
  SI metre-per-second readout of a speed magnitude
\end{definition}
\begin{proof}
  This is the declared model definition of speed In Meters Per Second.
\end{proof}
\begin{definition}[acceleration In Meters Per Second Squared]
  \label{def:physics:phyx-mini-0793:phyxminiproblems-problemphyxmini0793-accelerationinmeterspersecondsquared}
  \lean{PhyXMiniProblems.ProblemPhyXMini0793.accelerationInMetersPerSecondSquared}
  \uses{def:physics:phyx-mini-0793:phyxminiproblems-problemphyxmini0793-signedaccelerationquantity, def:physics:phyx-mini-0793:phyxminiproblems-problemphyxmini0793-accelerationreadout}
  SI metre-per-second-squared readout of a signed acceleration
\end{definition}
\begin{proof}
  This is the declared model definition of acceleration In Meters Per Second Squared.
\end{proof}
\begin{definition}[Cardinal Direction]
  \label{def:physics:phyx-mini-0793:phyxminiproblems-problemphyxmini0793-cardinaldirection}
  \lean{PhyXMiniProblems.ProblemPhyXMini0793.CardinalDirection}
  Cardinal directions used to interpret the horizontal axis and arrows
\end{definition}
\begin{proof}
  This is the declared model definition of Cardinal Direction.
\end{proof}
\begin{definition}[Figure Object]
  \label{def:physics:phyx-mini-0793:phyxminiproblems-problemphyxmini0793-figureobject}
  \lean{PhyXMiniProblems.ProblemPhyXMini0793.FigureObject}
  Distinct objects visible in the supplied raster
\end{definition}
\begin{proof}
  This is the declared model definition of Figure Object.
\end{proof}
\begin{definition}[Figure Label]
  \label{def:physics:phyx-mini-0793:phyxminiproblems-problemphyxmini0793-figurelabel}
  \lean{PhyXMiniProblems.ProblemPhyXMini0793.FigureLabel}
  Literal symbolic labels visible in the supplied raster
\end{definition}
\begin{proof}
  This is the declared model definition of Figure Label.
\end{proof}
\begin{definition}[Supplied Motorcycle Figure]
  \label{def:physics:phyx-mini-0793:phyxminiproblems-problemphyxmini0793-suppliedmotorcyclefigure}
  \lean{PhyXMiniProblems.ProblemPhyXMini0793.SuppliedMotorcycleFigure}
  \uses{def:physics:phyx-mini-0793:phyxminiproblems-problemphyxmini0793-signedpositionquantity, def:physics:phyx-mini-0793:phyxminiproblems-problemphyxmini0793-timequantity, def:physics:phyx-mini-0793:phyxminiproblems-problemphyxmini0793-signedvelocityquantity, def:physics:phyx-mini-0793:phyxminiproblems-problemphyxmini0793-signedaccelerationquantity, def:physics:phyx-mini-0793:phyxminiproblems-problemphyxmini0793-cardinaldirection, def:physics:phyx-mini-0793:phyxminiproblems-problemphyxmini0793-figureobject, def:physics:phyx-mini-0793:phyxminiproblems-problemphyxmini0793-figurelabel}
  Typed transcription of the primary bitmap. The later position and velocity remain independent dimensionful quantities even though the image marks their scalar values with question marks
\end{definition}
\begin{proof}
  This is the declared model definition of Supplied Motorcycle Figure.
\end{proof}
\begin{definition}[Motorcycle Motion Setup]
  \label{def:physics:phyx-mini-0793:phyxminiproblems-problemphyxmini0793-motorcyclemotionsetup}
  \lean{PhyXMiniProblems.ProblemPhyXMini0793.MotorcycleMotionSetup}
  \uses{def:physics:phyx-mini-0793:phyxminiproblems-problemphyxmini0793-signedpositionquantity, def:physics:phyx-mini-0793:phyxminiproblems-problemphyxmini0793-timequantity, def:physics:phyx-mini-0793:phyxminiproblems-problemphyxmini0793-signedvelocityquantity, def:physics:phyx-mini-0793:phyxminiproblems-problemphyxmini0793-speedquantity, def:physics:phyx-mini-0793:phyxminiproblems-problemphyxmini0793-signedaccelerationquantity, def:physics:phyx-mini-0793:phyxminiproblems-problemphyxmini0793-suppliedmotorcyclefigure}
  The independent physical trajectory and the event selected by the question. In particular, queryTime and the position at that time are not defined from the recorded answer
\end{definition}
\begin{proof}
  This is the declared model definition of Motorcycle Motion Setup.
\end{proof}
\begin{definition}[Matches Problem Statement]
  \label{def:physics:phyx-mini-0793:phyxminiproblems-problemphyxmini0793-matchesproblemstatement}
  \lean{PhyXMiniProblems.ProblemPhyXMini0793.MatchesProblemStatement}
  \uses{def:physics:phyx-mini-0793:phyxminiproblems-problemphyxmini0793-positioninmeters, def:physics:phyx-mini-0793:phyxminiproblems-problemphyxmini0793-timeinseconds, def:physics:phyx-mini-0793:phyxminiproblems-problemphyxmini0793-velocityinmeterspersecond, def:physics:phyx-mini-0793:phyxminiproblems-problemphyxmini0793-speedinmeterspersecond, def:physics:phyx-mini-0793:phyxminiproblems-problemphyxmini0793-accelerationinmeterspersecondsquared, def:physics:phyx-mini-0793:phyxminiproblems-problemphyxmini0793-motorcyclemotionsetup}
  Numerical data from the prose, including the condition selecting the event in the question. No queried position or answer label occurs here
\end{definition}
\begin{proof}
  This is the declared model definition of Matches Problem Statement.
\end{proof}
\begin{definition}[Matches Supplied Motorcycle Figure]
  \label{def:physics:phyx-mini-0793:phyxminiproblems-problemphyxmini0793-matchessuppliedmotorcyclefigure}
  \lean{PhyXMiniProblems.ProblemPhyXMini0793.MatchesSuppliedMotorcycleFigure}
  \uses{def:physics:phyx-mini-0793:phyxminiproblems-problemphyxmini0793-positioninmeters, def:physics:phyx-mini-0793:phyxminiproblems-problemphyxmini0793-timeinseconds, def:physics:phyx-mini-0793:phyxminiproblems-problemphyxmini0793-velocityinmeterspersecond, def:physics:phyx-mini-0793:phyxminiproblems-problemphyxmini0793-accelerationinmeterspersecondsquared, def:physics:phyx-mini-0793:phyxminiproblems-problemphyxmini0793-motorcyclemotionsetup}
  Literal objects, labels, orientations, and numerical annotations read from image 793.png. The 2.0 s image snapshot is tied to the same trajectory but is deliberately not equated with queryTime
\end{definition}
\begin{proof}
  This is the declared model definition of Matches Supplied Motorcycle Figure.
\end{proof}
\begin{definition}[Has Physical Eastward Query Event]
  \label{def:physics:phyx-mini-0793:phyxminiproblems-problemphyxmini0793-hasphysicaleastwardqueryevent}
  \lean{PhyXMiniProblems.ProblemPhyXMini0793.HasPhysicalEastwardQueryEvent}
  \uses{def:physics:phyx-mini-0793:phyxminiproblems-problemphyxmini0793-timeinseconds, def:physics:phyx-mini-0793:phyxminiproblems-problemphyxmini0793-velocityinmeterspersecond, def:physics:phyx-mini-0793:phyxminiproblems-problemphyxmini0793-accelerationinmeterspersecondsquared, def:physics:phyx-mini-0793:phyxminiproblems-problemphyxmini0793-motorcyclemotionsetup}
  Direction and temporal facts selecting the physical future event. These facts contain neither its position nor a displayed answer value
\end{definition}
\begin{proof}
  This is the declared model definition of Has Physical Eastward Query Event.
\end{proof}
\begin{definition}[Satisfies Constant Acceleration Kinematics]
  \label{def:physics:phyx-mini-0793:phyxminiproblems-problemphyxmini0793-satisfiesconstantaccelerationkinematics}
  \lean{PhyXMiniProblems.ProblemPhyXMini0793.SatisfiesConstantAccelerationKinematics}
  \uses{def:physics:phyx-mini-0793:phyxminiproblems-problemphyxmini0793-positionreadout, def:physics:phyx-mini-0793:phyxminiproblems-problemphyxmini0793-timereadout, def:physics:phyx-mini-0793:phyxminiproblems-problemphyxmini0793-velocityreadout, def:physics:phyx-mini-0793:phyxminiproblems-problemphyxmini0793-speedreadout, def:physics:phyx-mini-0793:phyxminiproblems-problemphyxmini0793-accelerationreadout, def:physics:phyx-mini-0793:phyxminiproblems-problemphyxmini0793-motorcyclemotionsetup}
  Standard one-dimensional constant-acceleration kinematics, stated in every coherent choice of units and relative to the initial event. This interface contains no value for the queried position, no 55 m, and no answer label
\end{definition}
\begin{proof}
  This is the declared model definition of Satisfies Constant Acceleration Kinematics.
\end{proof}
\begin{definition}[query Elapsed Time Seconds]
  \label{def:physics:phyx-mini-0793:phyxminiproblems-problemphyxmini0793-queryelapsedtimeseconds}
  \lean{PhyXMiniProblems.ProblemPhyXMini0793.queryElapsedTimeSeconds}
  \uses{def:physics:phyx-mini-0793:phyxminiproblems-problemphyxmini0793-timeinseconds, def:physics:phyx-mini-0793:phyxminiproblems-problemphyxmini0793-motorcyclemotionsetup}
  Elapsed SI time from the stated initial event to the queried event
\end{definition}
\begin{proof}
  This is the declared model definition of query Elapsed Time Seconds.
\end{proof}
\begin{definition}[queried Position East Of Signpost Meters]
  \label{def:physics:phyx-mini-0793:phyxminiproblems-problemphyxmini0793-queriedpositioneastofsignpostmeters}
  \lean{PhyXMiniProblems.ProblemPhyXMini0793.queriedPositionEastOfSignpostMeters}
  \uses{def:physics:phyx-mini-0793:phyxminiproblems-problemphyxmini0793-positioninmeters, def:physics:phyx-mini-0793:phyxminiproblems-problemphyxmini0793-motorcyclemotionsetup}
  Requested signed position relative to the city-limits signpost, in metres
\end{definition}
\begin{proof}
  This is the declared model definition of queried Position East Of Signpost Meters.
\end{proof}
\begin{lemma}[query Elapsed Time from speed]
  \label{lem:physics:phyx-mini-0793:phyxminiproblems-problemphyxmini0793-queryelapsedtime-from-speed}
  \lean{PhyXMiniProblems.ProblemPhyXMini0793.queryElapsedTime_from_speed}
  \uses{def:physics:phyx-mini-0793:phyxminiproblems-problemphyxmini0793-motorcyclemotionsetup, def:physics:phyx-mini-0793:phyxminiproblems-problemphyxmini0793-matchesproblemstatement, def:physics:phyx-mini-0793:phyxminiproblems-problemphyxmini0793-hasphysicaleastwardqueryevent, def:physics:phyx-mini-0793:phyxminiproblems-problemphyxmini0793-satisfiesconstantaccelerationkinematics, def:physics:phyx-mini-0793:phyxminiproblems-problemphyxmini0793-queryelapsedtimeseconds}
  The 25 m/s event occurs 2.5 s after the stated initial event
\end{lemma}
\begin{proof}
  The conclusion follows from the governing relations and figure readouts named in the statement of query Elapsed Time from speed.
\end{proof}
\begin{lemma}[queried Position from kinematics]
  \label{lem:physics:phyx-mini-0793:phyxminiproblems-problemphyxmini0793-queriedposition-from-kinematics}
  \lean{PhyXMiniProblems.ProblemPhyXMini0793.queriedPosition_from_kinematics}
  \uses{def:physics:phyx-mini-0793:phyxminiproblems-problemphyxmini0793-motorcyclemotionsetup, def:physics:phyx-mini-0793:phyxminiproblems-problemphyxmini0793-matchesproblemstatement, def:physics:phyx-mini-0793:phyxminiproblems-problemphyxmini0793-hasphysicaleastwardqueryevent, def:physics:phyx-mini-0793:phyxminiproblems-problemphyxmini0793-satisfiesconstantaccelerationkinematics, def:physics:phyx-mini-0793:phyxminiproblems-problemphyxmini0793-queriedpositioneastofsignpostmeters}
  Constant-acceleration kinematics gives the requested 55 m position
\end{lemma}
\begin{proof}
  The conclusion follows from the governing relations and figure readouts named in the statement of queried Position from kinematics.
\end{proof}
\begin{definition}[Answer Choice]
  \label{def:physics:phyx-mini-0793:phyxminiproblems-problemphyxmini0793-answerchoice}
  \lean{PhyXMiniProblems.ProblemPhyXMini0793.AnswerChoice}
  Labels of the four positions printed in the answer list
\end{definition}
\begin{proof}
  This is the declared model definition of Answer Choice.
\end{proof}
\begin{definition}[displayed Position Meters]
  \label{def:physics:phyx-mini-0793:phyxminiproblems-problemphyxmini0793-displayedpositionmeters}
  \lean{PhyXMiniProblems.ProblemPhyXMini0793.displayedPositionMeters}
  \uses{def:physics:phyx-mini-0793:phyxminiproblems-problemphyxmini0793-answerchoice}
  Position in metres printed beside each answer label
\end{definition}
\begin{proof}
  This is the declared model definition of displayed Position Meters.
\end{proof}
\begin{definition}[recorded Answer Choice]
  \label{def:physics:phyx-mini-0793:phyxminiproblems-problemphyxmini0793-recordedanswerchoice}
  \lean{PhyXMiniProblems.ProblemPhyXMini0793.recordedAnswerChoice}
  \uses{def:physics:phyx-mini-0793:phyxminiproblems-problemphyxmini0793-answerchoice}
  Answer label recorded by the supplied dataset
\end{definition}
\begin{proof}
  This is the declared model definition of recorded Answer Choice.
\end{proof}
\begin{definition}[Matches Displayed Position]
  \label{def:physics:phyx-mini-0793:phyxminiproblems-problemphyxmini0793-matchesdisplayedposition}
  \lean{PhyXMiniProblems.ProblemPhyXMini0793.MatchesDisplayedPosition}
  \uses{def:physics:phyx-mini-0793:phyxminiproblems-problemphyxmini0793-motorcyclemotionsetup, def:physics:phyx-mini-0793:phyxminiproblems-problemphyxmini0793-queriedpositioneastofsignpostmeters, def:physics:phyx-mini-0793:phyxminiproblems-problemphyxmini0793-answerchoice, def:physics:phyx-mini-0793:phyxminiproblems-problemphyxmini0793-displayedpositionmeters}
  A displayed position agrees with the position of the queried event
\end{definition}
\begin{proof}
  This is the declared model definition of Matches Displayed Position.
\end{proof}
\begin{definition}[Is Unique Matching Displayed Choice]
  \label{def:physics:phyx-mini-0793:phyxminiproblems-problemphyxmini0793-isuniquematchingdisplayedchoice}
  \lean{PhyXMiniProblems.ProblemPhyXMini0793.IsUniqueMatchingDisplayedChoice}
  \uses{def:physics:phyx-mini-0793:phyxminiproblems-problemphyxmini0793-motorcyclemotionsetup, def:physics:phyx-mini-0793:phyxminiproblems-problemphyxmini0793-answerchoice, def:physics:phyx-mini-0793:phyxminiproblems-problemphyxmini0793-matchesdisplayedposition}
  Exactly one displayed label agrees with the requested physical position
\end{definition}
\begin{proof}
  This is the declared model definition of Is Unique Matching Displayed Choice.
\end{proof}
% --- Archon declaration topology end ---
% --- Archon physics formalization source end ---
